\skyrius{Produkto kūrimas}
\subsection{Kūrimo eiga ir esminiai sprendimai}
\gaires{Atrinkite svarbiausius kūrimo etapus, o ne rašykite kasdienį veiksmų
sąrašą. Kiekviename etape paaiškinkite: ko siekėte, ką padarėte, kas nepavyko,
kokį pakeitimą pasirinkote ir kuo jis padėjo. Nurodykite naudotus metodus,
priemones bei svarbias versijas.}
\pildyti{Pirmojo prototipo, jo tobulinimo ir galutinės versijos aprašymas.
Pateikite pasirinkimus paaiškinančių įrodymų: brėžinių, ekrano vaizdų,
bandymo įrašų, versijų žymų ar trumpų kodo fragmentų.}

\subsection{Tarpiniai rezultatai ir plano pakeitimai}
\pildyti{Kokį tarpinį rezultatą ir kada aptarėte su vadovu? Kokią pastabą
gavote ir kaip į ją atsižvelgėte? Kokie reikalavimai ar plano terminai pasikeitė,
kodėl ir kaip tai paveikė rezultatą? Remkitės \ref{app:zurnalas} priedu.}

\subsection{Asmeninis indėlis ir perimti komponentai}
\gaires{Aiškiai atskirkite savo darbą nuo bibliotekų, pradinių projektų,
duomenų rinkinių, kitų žmonių pagalbos ir sugeneruoto turinio. Atviro kodo
bibliotekos naudojimas savaime neparodo jūsų indėlio: svarbu, ką suprojektavote,
sukūrėte ir patikrinote patys.}
\pildyti{Ką sukūrėte savarankiškai? Ką panaudojote ar pritaikėte ir kokiomis
licencijos sąlygomis? Jei naudotasi DI priemonėmis, nurodykite priemonę,
naudojimo paskirtį, paveiktą darbo dalį ir savo atliktą patikrą.
Laikykitės mokykloje suderintų DI naudojimo taisyklių.}

\pildyti{Jei dirbote grupėje, tiksliai įvardykite kiekvieno autoriaus atliktas
dalis ir savo indėlį. Jei projekte dalyvavo naudotojai, paaiškinkite jų
informavimą, duomenų apsaugą ir kaip išvengėte nereikalingų asmens duomenų rinkimo.}
